\documentclass[10pt,twocolumn]{article}

\usepackage[margin=1.9cm,top=2.2cm,bottom=2.2cm]{geometry}
\usepackage{amsmath}
\usepackage{amssymb}
\usepackage{hyperref}
\usepackage{times}
\usepackage{titlesec}
\usepackage{booktabs}
\usepackage{microtype}
\usepackage{parskip}
\hypersetup{colorlinks=true, linkcolor=blue, citecolor=blue, urlcolor=blue}

\titleformat{\section}{\normalfont\bfseries\small\centering}{}{0em}{}
\titlespacing{\section}{0pt}{6pt}{3pt}

\setlength{\columnsep}{0.5cm}
\setlength{\parskip}{3pt}
\setlength{\parindent}{1em}

\usepackage{fancyhdr}
\pagestyle{fancy}
\fancyhf{}
\fancyhead[L]{\small\itshape BCT Superfluid Lattice Model --- Letter 30}
\fancyhead[R]{\small\itshape March 2026}
\fancyfoot[C]{\thepage}
\renewcommand{\headrulewidth}{0.4pt}

\begin{document}

\twocolumn[{%
\begin{center}
{\large\bfseries The Neutron--Proton Mass Difference at Sub-0.1\%:}\\[3pt]
{\large\bfseries BCT Colour Casimir SU(6) Theorem}\\[8pt]
{\normalsize Michel Robert Cabri\'{e}$^{1,\,*}$}\\[2pt]
{\small $^1$\textit{Independent Researcher;  ORCID: \href{https://orcid.org/0009-0007-9561-9859}{0009-0007-9561-9859}}}\\[2pt]
{\small (Dated: March 2026)\quad BCT Letter 30}\\[6pt]
\end{center}
\begin{quotation}
\small\noindent
We apply the BCT Colour Casimir Theorem established in Letter~29 to the neutron--proton mass
difference $\Delta m_N = m_n - m_p$. At leading order, the BCT SU(6) spin-flavour theorem gives
$\Delta m_N^{\rm LO} = \Lambda_{\rm QCD}\,\alpha_0\,\pi/4 = 1.280$~MeV (Phase~48B, $-1.03\%$).
Applying the NLO Casimir correction $F_{\rm SU(6)} \to F_{\rm SU(6)}(1+C_F\alpha_0)$ with
$C_F = (N_c^2-1)/(2N_c) = 4/3$ yields
\[
  \Delta m_N = \frac{\Lambda\alpha_0\pi}{4}(1+C_F\alpha_0) = 1.2927\,\text{MeV}\,,\quad
  \delta = -0.052\%\,,
\]
against the observed value $1.29333$~MeV. This is \textbf{Prediction~\#105} of the BCT programme
and the fourth sub-0.1\% result obtained in Phase~97. Combined with Letter~29
(Predictions~\#102--104), this session delivers four simultaneous sub-0.1\% nuclear predictions
--- kaon mass, proton mass, neutron mass, and neutron--proton mass difference --- from a single
NLO colour correction with zero new free parameters. The $N$-class coefficient $C_F = 4/3$
appears identically in all four predictions, providing strong internal consistency for the BCT
NLO colour structure.
\end{quotation}
\vspace{6pt}
}]

\section{Introduction}

The neutron--proton mass difference $\Delta m_N = m_n - m_p = 1.29333$~MeV is one of the most
cosmologically critical numbers in physics. Were it negative, the proton would be unstable; were
it much larger than the electron mass, Big Bang nucleosynthesis would produce almost no helium.
Its precise value controls the $n/p$ freeze-out ratio and the primordial $^4$He abundance
$Y_p \approx 0.25$. Deriving it from first principles is therefore a sharp test of any
fundamental theory.

The BCT programme has addressed $\Delta m_N$ in two previous phases. Phase~48B~[1] established
the leading-order result $\Delta m_N^{\rm LO} = \Lambda\alpha_0\pi/4 = 1.280$~MeV ($-1.03\%$)
from the BCT SU(6) spin-flavour theorem. Phase~81D~[1] improved this to $1.285$~MeV
($-0.647\%$) by separating the strong NLO and electromagnetic Cottingham contributions. Both are
sub-1\% but not sub-0.1\%.

Letter~29~[3] introduced the BCT Colour Casimir Theorem: when meson modes in a hadronic
decomposition carry their own NLO colour corrections, the baryon NLO coefficient upgrades from
$N = 3/4$ to $C_F = (N_c^2-1)/(2N_c) = 4/3$, the SU(3) colour Casimir. This theorem produced
Predictions~\#102--104: the kaon, proton, and neutron masses all sub-0.1\%. Here we apply the
same theorem to the nucleon mass splitting and obtain Prediction~\#105.

\section{The BCT SU(6) Theorem at LO}

The proton (u,u,d) and neutron (u,d,d) differ by a single $u \leftrightarrow d$ quark exchange.
The BCT $T_d$ void lattice assigns $u$-type quarks to octahedral voids (radius $r_{\rm oct}$)
and $d$-type quarks to tetrahedral voids (radius $r_{\rm tet}$). The $D_{4h}$ tetragonal
distortion of the BCT lattice ($c/a = \sqrt{2}$) breaks $u$-$d$ isospin symmetry, generating
the quark mass difference~[1]:
\begin{equation}
m_d - m_u = \Lambda\alpha_0\pi/2\,,
\label{eq:quarkdiff}
\end{equation}
with all quantities from the three BCT inputs $\{r_{\rm oct}, r_{\rm tet}, \Lambda_{\rm QCD}\}$.

The BCT SU(6) spin-flavour overlap integral for the single-quark substitution in the
non-relativistic $T_d$ nucleon wave function gives $F_{\rm SU(6)} = 1/2$~[1]. Physically, each
nucleon contains three quarks, but the $T_d$ doublet geometry of the $(u,d)$ pair suppresses
the single-quark matrix element by $1/2$ rather than the naive $1/3$. The LO prediction is:
\begin{equation}
\Delta m_N^{\rm LO} = (m_d - m_u)\cdot F_{\rm SU(6)}
  = \frac{\Lambda\alpha_0\pi}{4} = 1.280020\,\text{MeV}\,,
\label{eq:DmLO}
\end{equation}
with observed $\Delta m_N = 1.29333$~MeV and error $-1.03\%$.

\section{NLO Casimir Correction --- Prediction \#105}

\textit{The Casimir upgrade.} The BCT Colour Casimir Theorem (Letter~29) states that when
hadronic states involve NLO colour corrections, the baryon self-energy coefficient upgrades to
the SU(3) Casimir $C_F$. For the nucleon mass splitting, the relevant correction is to the
SU(6) factor $F_{\rm SU(6)}$: the single quark being exchanged ($u \to d$) carries colour
charge $C_F$ and acquires an NLO colour-loop dressing. The upgrade is:
\begin{equation}
F_{\rm SU(6)} = \frac{1}{2}
  \;\longrightarrow\;
  F_{\rm SU(6)}^{\rm NLO} = \frac{1}{2}(1+C_F\alpha_0)\,,
\label{eq:FupgradE}
\end{equation}
where $C_F = (N_c^2-1)/(2N_c) = 4/3$ and $N_c = 3$ is derived from the BCT $D_4$ triality
theorem~[1].

\textit{The NLO formula.} Substituting Eq.~(\ref{eq:FupgradE}) into Eq.~(\ref{eq:DmLO}):
\begin{equation}
\Delta m_N^{\rm NLO} = \frac{\Lambda\alpha_0\pi}{4}(1+C_F\alpha_0)\,.
\label{eq:DmNLO}
\end{equation}

\textit{Numerical evaluation.} With $\Lambda = 220$~MeV, $\alpha_0 = r_{\rm oct}\,r_{\rm tet}/\pi
= 0.00740806$, $C_F = 4/3$:
\begin{align}
\Lambda\alpha_0\pi/4 &= 1.280020\,\text{MeV}, \label{eq:n1}\\
C_F\alpha_0 &= \tfrac{4}{3}\times 0.00740806 = 0.00987741, \label{eq:n2}\\
1 + C_F\alpha_0 &= 1.00987741, \label{eq:n3}
\end{align}
\begin{equation}
\Delta m_N^{\rm NLO} = 1.292663\,\text{MeV}\,,\quad
\delta = -0.052\%\;\;(\star\,\text{sub-0.1\%}).
\label{eq:result}
\end{equation}
against the CODATA value $\Delta m_N^{\rm obs} = 1.29333$~MeV. This is \textbf{Prediction~\#105}.

\section{The $N$-Class Quartet}

The BCT Colour Casimir Theorem now governs four simultaneous predictions in the nuclear sector,
all from Phase~97:

\begin{table}[h]
\centering\small
\caption{Phase~97 nuclear sub-0.1\% quartet. All from
$\{r_{\rm oct}, r_{\rm tet}, \Lambda_{\rm QCD}\}$; $C_F = 4/3$ throughout.}
\label{tab:quartet}
\begin{tabular}{@{}llccc@{}}
\toprule
Pred. & Observable & BCT Value & Obs. & Error \\
\midrule
\#102 & $m_K$ (NLO)  & 493.534~MeV & 493.677 & $-0.029\%\star$ \\
\#103 & $m_p$ (NLO)  & 938.288~MeV & 938.272 & $+0.002\%\star$ \\
\#104 & $m_n$ (NLO)  & 939.541~MeV & 939.565 & $-0.003\%\star$ \\
\#105 & $\Delta m_N$ & 1.2927~MeV  & 1.2933  & $-0.052\%\star$ \\
\bottomrule
\multicolumn{5}{@{}l@{}}{\footnotesize $\star$ Sub-0.1\%. All $N$-class: $N = C_F = 4/3$.}
\end{tabular}
\end{table}

The uniformity of the $N$-class coefficient $C_F = 4/3$ across all four observables is a
non-trivial internal consistency check. In each case the physical mechanism is the same: a
quark undergoing an NLO colour-loop interaction acquires a self-energy proportional to
$C_F\alpha_0$, where $C_F$ is the quadratic Casimir of the fundamental SU(3) representation.

\textit{Comparison with Phase~81D.} The Phase~81D result~[1] obtained $\Delta m_N = 1.2846$~MeV
($-0.647\%$) by separating the contribution into a strong NLO term $2m_\pi\alpha_0(1+3\alpha_0)
= 2.045$~MeV and an electromagnetic Cottingham term $-0.76$~MeV (taken from standard QCD+QED
calculations). The present result achieves better precision without invoking the Cottingham
separation, using instead the compact formula Eq.~(\ref{eq:DmNLO}) which treats the full
$\Delta m_N$ holistically through the BCT geometric approach.

\textit{Cosmological note.} The BCT formula Eq.~(\ref{eq:DmNLO}) evaluates to
$\Delta m_N/m_N = \alpha_0\pi C_F\alpha_0/4\ldots$, a tiny but precise fraction of the nucleon
mass set entirely by void geometry. BCT naturally generates $\Delta m_N$ in the window
$[m_e, 3\,\text{MeV}]$ required for both proton stability and Big Bang nucleosynthesis ---
without any parameter tuning.

\section{Conclusion}

Applying the BCT Colour Casimir Theorem to the SU(6) spin-flavour factor governing the
neutron--proton mass difference gives:
\begin{align}
\Delta m_N &= \frac{\Lambda\alpha_0\pi}{4}(1+C_F\alpha_0)
            = 1.2927\,\text{MeV}\,,\notag\\
           &\quad\delta = -0.052\%\,,
\label{eq:concl}
\end{align}
from $\{r_{\rm oct}, r_{\rm tet}, \Lambda_{\rm QCD}\}$ alone, with no new parameters.
Combined with Letter~29, Phase~97 delivers four simultaneous sub-0.1\% predictions
(\#102--\#105) for the four most fundamental nuclear observables --- kaon mass, proton mass,
neutron mass, and nucleon mass splitting --- all governed by the single Casimir coefficient
$C_F = (N_c^2-1)/(2N_c) = 4/3$ derived from BCT $D_4$ triality.

The BCT programme stands at \textbf{105 sub-1\% predictions} and \textbf{30 sub-0.1\%
predictions}, all from $\{r_{\rm oct}, r_{\rm tet}, \Lambda_{\rm QCD}\}$.

\smallskip
\noindent\textit{Data availability.} All results reproduced from the formulas given; no datasets
beyond CODATA/PDG values cited are used.

\bigskip
\noindent$^*$ \href{mailto:ZeroFreeParameters@gmail.com}{\texttt{ZeroFreeParameters@gmail.com}}

\begin{thebibliography}{4}

\bibitem{BCT_Monograph}
M.~R. Cabri\'{e}, \textit{BCT Superfluid Lattice Model: Complete Derivation of Standard Model
Parameters}, Zenodo (2026),
\href{https://doi.org/10.5281/zenodo.18884976}{doi:10.5281/zenodo.18884976}.

\bibitem{BCT_PRL}
M.~R. Cabri\'{e}, \textit{Zero Free Parameters: Deriving 92 Standard Model Observables from
BCT Vacuum Geometry}, Zenodo (2026),
\href{https://doi.org/10.5281/zenodo.18884415}{doi:10.5281/zenodo.18884415}.

\bibitem{BCT_L29}
M.~R. Cabri\'{e}, BCT Letter~29: The NLO Colour Correction to the Strange Sector ---
Three Sub-0.1\% Predictions, Zenodo (2026),
\href{https://doi.org/10.5281/zenodo.18884415}{doi:10.5281/zenodo.18884415}.

\bibitem{PDG2022}
R.~L. Workman et al.\ (Particle Data Group), Prog.\ Theor.\ Exp.\ Phys.\ \textbf{2022},
083C01 (2022).

\end{thebibliography}

\end{document}
